\chapter{Introduction}

\section{Motivation}

Continuous-time generative models learn to transform a simple noise distribution
into a complex data distribution by following a learned trajectory in pixel space.
At inference, this trajectory is integrated numerically: the model is called once
per integration step, and each call is known as a \emph{function evaluation}
(NFE). Reducing NFE without sacrificing image quality is therefore a core research
objective, since it directly controls generation latency and compute cost.

Flow Matching (FM)~\cite{lipman2023flow} addresses this by training a neural
network to predict the instantaneous velocity of a prescribed interpolation path
between data and noise, without simulating the resulting ODE during training. This
conditional-path approach yields stable gradients and fast training. However,
standard FM still requires many NFE for accurate integration, because the learned
field can vary significantly along the trajectory.

Mean Flow (MF)~\cite{geng2025mean} takes a different approach: instead of
predicting a local velocity at a single time, the network predicts the
\emph{average velocity} over an interval $[r, t]$. When the interval spans the
entire noise-to-data range $[0, 1]$, a single network call directly displaces a
noise sample to a data sample---making one-step generation theoretically exact
rather than an approximation.

A third configuration tests logit-normal time sampling, motivated by
rectified-flow work~\cite{esser2024scaling} showing that concentrating training
samples near intermediate noise levels improves generation quality under the
same number of training steps.

\section{Problem Statement}

Algorithm comparisons become unreliable when the model architecture,
preprocessing pipeline, evaluation metrics, random seed, or log interpretation
differ between runs. This project therefore combines implementation with evidence
control: all three algorithms share a single backbone, data pipeline, evaluator,
and JSONL logging schema, while every remaining configuration difference is
explicitly disclosed and its effect on comparability is assessed.

\section{Aim, Objectives, and Research Questions}

The aim is to build and critically evaluate a reproducible educational system for
unconditional CIFAR-10 image generation using three flow-based generative
objectives. Specific objectives are:

\begin{enumerate}
  \item Implement FM, FM-LN, and MF as interchangeable wrappers around a shared
        backbone and training engine.
  \item Preserve identical evaluation infrastructure across all algorithms: same
        FID reference, same NFE grid, same JSONL schema.
  \item Measure and record training loss, sampling time, GPU memory, FID, and IS
        across NFE $\in\{1,5,10,20\}$.
  \item Provide a local browser-based demonstration of training history playback
        and live inference.
  \item Audit all JSONL logs to detect and resolve duplicate epoch entries before
        reporting.
\end{enumerate}

The project is guided by the following research questions:

\begin{description}
  \item[RQ1.] How do the three training objectives behave during training---in
              terms of loss trajectory, convergence, and stability?
  \item[RQ2.] How does sample quality (FID, IS) change with NFE for each
              algorithm?
  \item[RQ3.] What quality--latency trade-offs are observed at matched and
              mismatched compute budgets?
  \item[RQ4.] Does this MF implementation realise competitive one-step generation
              relative to multi-step FM?
  \item[RQ5.] Can every reported claim be traced back to a configuration file,
              checkpoint, JSONL log, and UI demonstration?
\end{description}

\section{Contributions}

The deliverable comprises:

\begin{itemize}
  \item Complete FM, FM-LN, and MF implementations sharing one 6.35M-parameter
        SimpleUNet backbone.
  \item Ten-epoch checkpoint archives for preserving submitted model states.
  \item Structured JSONL evidence with a documented audit protocol for handling
        duplicate restart rows.
  \item Cached real-image FID reference statistics with metadata validation to
        prevent silent distribution mismatch.
  \item FID and IS evaluation via TorchMetrics, batched to avoid GPU overflow.
  \item A browser inference API (\texttt{inference\_server.py}) and a training
        history dashboard that synchronises checkpoint images with loss curves.
  \item A deterministic duplicate-log audit script that validates epoch
        contiguity, parameter invariance, and NFE grid completeness.
  \item This report, with all figures generated from logged artefacts and all
        tables cross-checked against source JSONL.
\end{itemize}
